\documentclass[10pt,a4paper]{amsart}
\usepackage[margin=19mm]{geometry}
\usepackage{amsmath,amssymb,mathtools}
\usepackage[scr=rsfs,cal=euler]{mathalfa}
\usepackage{xfrac}
\usepackage[unicode,hidelinks,bookmarks=false]{hyperref}
\newcommand{\ie}{i.e.\ }
\newcommand{\cf}{cf.\ }
\newcommand{\resp}{respectively\ }
\newcommand{\Subsection}{\subsection}
\providecommand{\nobreakdash}{\nobreak\hbox{-}}
\setlength{\emergencystretch}{2em}
\multlinegap0pt
\usepackage{bm}
\usepackage{bbm}
\usepackage{dsfont}
\usepackage{xspace}
\usepackage{cancel}
\usepackage{mathtools}
\usepackage{amsthm}
\makeatletter
\@ifundefined{theorem}{\newtheorem{theorem}{Theorem}}{}
\@ifundefined{corollary}{\newtheorem{corollary}[theorem]{Corollary}}{}
\@ifundefined{lemma}{\newtheorem{lemma}[theorem]{Lemma}}{}
\@ifundefined{proposition}{\newtheorem{proposition}[theorem]{Proposition}}{}
\makeatother
\makeatletter
\def\Cl{{\rm Cl}}
\def\Q{{\rm \mathbb{Q}}}
\def\Sel{{\rm Sel}}
\def\Z{{\rm \mathbb{Z}}}
\def\avg{{\rm avg}}
\expandafter\ifx\csname derivation\endcsname\relax
\newenvironment{derivation}{%
  \renewcommand{\proofname}{Derivation}\proof}{\endproof}
\fi
\expandafter\ifx\csname exercise\endcsname\relax
\newenvironment{exercise}[1][Exercise]{\begin{trivlist}
\item[\hskip \labelsep {\bfseries #1}]}{\end{trivlist}}
\fi
\def\un{{\rm un}}
\makeatother
\title[Hecke reciprocity and class groups]{Hecke reciprocity and class groups}
\author{Ari Shnidman, Artane Siad}
\date{Source: arXiv:2506.13749v2 (CC BY 4.0)}
\begin{document}
\maketitle

\noindent\emph{Source and licence.} Hecke reciprocity and class groups. Version arXiv:2506.13749v2 (CC BY 4.0), distributed under the Creative Commons Attribution 4.0 International licence, as recorded in the arXiv licence metadata for this exact version and shown on the abstract page https://arxiv.org/abs/2506.13749v2. Full rights notice: licenses/arxiv-2506.13749-CC-BY-4.0.txt.\par\smallskip

\noindent\textit{Editorial scope.} Counted unit: the Theorem whose source label is \texttt{thm: tame selmer avg} in the article (source0.tex lines 1052--1054, source label \texttt{thm: tame selmer avg}), delivered together with the article's own complete original proof of it (source0.tex lines 1056--1074, one \texttt{proof} environment of 19 lines). No mathematics is written, repaired or re-derived: statement and proof are the article's own text line for line. Required context retained uncounted: lem: local selmer ratio (lines 514--516); thm: wild selmer avg (lines 1007--1009); prop: tame failure II (lines 1026--1028); cor: Z bijection (lines 1046--1048). Every definition, display and label of those lines is retained unchanged. Omitted from this package and not redistributed: 1--513, 517--1006, 1010--1025, 1029--1045, 1049--1051, 1055--1055, 1075--1649. Nothing inside a retained excerpt is omitted or altered: every retained body is delivered exactly as the article's lines are written, so no omission receipt of any kind accompanies this package. Citations: the retained excerpts carry no citation command at all, so no bibliography entry of the article is transcribed and no reference is redistributed. Reference adaptation: none was needed and none was made; every reference inside the delivered unit resolves inside it, so no unresolved reference or citation is produced. Printed slips kept literal: no printed word, symbol, hypothesis or number of the counted statement or of its proof was altered. Layout and class: the article's own class is \texttt{amsart} with packages including amsfonts, amsmath, amscd, amsthm, fullpage, amssymb, centernot, enumerate, pb-diagram, mathrsfs, fontenc, seqsplit, tikz, tikz-cd; the wrapper is the lane's pinned \texttt{amsart} editorial wrapper at 10pt on A4, which restates the theorem-like environments the retained text needs and adds the article's own macro definitions that the retained text needs (\texttt{\textbackslash Cl}, \texttt{\textbackslash Q}, \texttt{\textbackslash Sel}, \texttt{\textbackslash Z}, \texttt{\textbackslash avg}, \texttt{\textbackslash un}), with extra wrapper packages (bm, bbm, dsfont, xspace, cancel, mathtools). The delivered numbering is the unit's own. Assets: no retained span contains a figure environment, a graphics-inclusion command, a PDF-page inclusion, a table or a TikZ picture; no image file is copied into this package, the package declares no asset dependency and no image file is redistributed with it. Licence: version arXiv:2506.13749v2 of this article is distributed under the Creative Commons Attribution 4.0 International licence, as recorded in the arXiv licence metadata for this exact version and shown on the abstract page https://arxiv.org/abs/2506.13749v2. A full rights notice is delivered at licenses/arxiv-2506.13749-CC-BY-4.0.txt. Characters: every retained excerpt is pure ASCII, so the delivered engine, XeLaTeX with its default Latin Modern text font, reports no missing character.
\begin{lemma}\label{lem: local selmer ratio}
    If $\ell \nmid [F: K]$ and $v$ is a finite prime of $K$, then  $|V_v| = |M(K_v)| = \ell^{N-1}$, where $N$ is the number of primes of $F$ above $v$. 
\end{lemma}

\begin{theorem}\label{thm: wild selmer avg}
    Ordering the wild fields $F_n = \Q(\sqrt[3]{n})$ by $|n|$, we have $\avg_n\, |\Cl_{F_n}[2]| = 2$.
\end{theorem}

\begin{proposition}\label{prop: tame failure II}
Suppose $n \equiv \pm 1\pmod{9}$. Then the non-trivial element of $V_{n,3} \subset H^1(\Q_3, M_n)$ is not contained in $\ker(q_{n,3})$. 
\end{proposition}

\begin{corollary}\label{cor: Z bijection}
For each $n \geq 1$, there exists a bijection $\Sel_2^\un(F_n) \to G(\Q)\backslash Y(\Q)^\un_n$. Moreover, there exists $N \geq 1$,  independent of $n$, such that each $G(\Q)$-orbit in $Y(\Q)^\un_n$ is represented by an element in $\frac{1}{N}V(\Z)$. 
\end{corollary}

\begin{theorem}\label{thm: tame selmer avg}
    Ordering the tame fields $F_n = \Q(\sqrt[3]{n})$ by $|n|$, we have $\avg_n\, |\Cl_{F_n}[2]|= 1 + \frac12 $.
\end{theorem}

\begin{proof}
 Let $U \subset \Z$ be the integers congruent to $\pm 1\pmod9$, and for each prime $p$, let $U_p$ be the closure of $U$ in $\Q_p$. Proceeding as in the proof of Theorem \ref{thm: wild selmer avg}, and using Corollary \ref{cor: Z bijection}, we obtain
\[\avg_d \#\Sel_2^\un(F_n) = 1 + 2\prod_p \nu_p,\] 
where this time 
\[\nu_p = \int_{n \in U_p} \frac{\#V_{n,p} \cap \ker(q_{n,p})}{\#M_n(\Q_p)}dn.\]
The computation for $p \neq 3$ is as before. For $p = 3$ we have:

\begin{lemma}
$\nu_3 = \frac12$.
\end{lemma}
\begin{proof}
Since $F_n$ is tame, the prime $3$ has splitting type $(1^21)$ and hence $\#V_{n,3} = \#M_n(\Q_3) = 2$ by Lemma \ref{lem: local selmer ratio}. On the other hand, we have $\#V_{n,3} \cap \ker(q_{n,3})= 1$ by Proposition \ref{prop: tame failure II}. 
\end{proof}


Putting this together, we obtain,
\[\avg_n \#\Sel_2^\un(F_n) = 1 + 2\prod_p \nu_p =1 + 2\cdot \frac12 \cdot \frac12 \cdot \prod_{p \nmid 3\infty}  1 =  1 + \frac12,\]
as desired.
\end{proof}

\end{document}
